%% 06-dashboards.tex — Dashboard links and how to read maps

\section{Dashboards and Maps}

\subsection*{The Static Dashboard}

The portfolio ships with a self-contained HTML dashboard at
\texttt{web/dashboard.html}.
Open it in any modern browser---no server, no login, no software installation.
The dashboard shows:

\begin{itemize}
  \item A map of every state's algorithmic districts for 2000, 2010, and 2020.
  \item Compactness scores for each district, color-coded from compact
        (dark green) to elongated (light yellow).
  \item Population balance: every district is within $\pm$0.5\% of the
        state's ideal district population.
  \item Side-by-side comparison with enacted 2020 districts where available.
\end{itemize}

\subsection*{How to Read a Bisection Map}

Each district in the map corresponds to one leaf of the bisection tree.
Districts that were separated early in the process (at the first or second
cut) tend to be large geographic regions: a northern half, a southern half.
Districts separated late are more local and more finely shaped.

Colors in the standard output represent the bisection \emph{round} at which
a district was finalized:
\begin{itemize}
  \item Round 1 (two pieces): wide, regional districts.
  \item Round 2 (four pieces): quadrant-level districts.
  \item Round 3 and later: neighborhood-scale districts in high-density states.
\end{itemize}

\subsection*{What the Round-by-Round Images Show}

The pipeline generates one image per round per state.
For a state with $k$ seats, there are $\lceil \log_2 k \rceil$ rounds.
Each image shows the current partition, with district boundaries
darkening as more cuts are made.

Reading the rounds from first to last, you see the algorithm
``zoom in'' from large geographic halves to individual districts.
This visualization makes it easy to verify that no round
produced a non-contiguous piece or a population imbalance.

\subsection*{Reproducing Any Figure}

Every figure in the portfolio can be regenerated from the public
Census data and the open-source \texttt{redist} CLI:

\begin{verbatim}
redist fetch --year 2020
redist state --state NC --year 2020 --version v1
redist map   --state NC --year 2020 --version v1
\end{verbatim}

\noindent The resulting maps and CSV data are placed in
\texttt{outputs/v1/2020/north\_carolina/}.
The full replication procedure is in Paper A.4.
